\section*{الفصل \ref{ch:b1:electrophilic-additions} --- الألكينات والألكاينات: الإضافات المحبة للإلكترونات}

\begin{solution}{exo:b1:electrophilic-additions:1}
2-كلوروبروبان؛ 2-برومو-2-ميثيل البروبان؛ 1-برومو-1-ميثيل الهكسان الحلقي؛
2-كلوروبوت-2-ين.
\end{solution}

\begin{solution}{exo:b1:electrophilic-additions:2}
1,2-ثنائي برومو الإيثان؛ 1,2-ثنائي برومو الهكسان الحلقي-\textit{trans}؛
1,2-ثنائي برومو البروبان.
\end{solution}

\begin{solution}{exo:b1:electrophilic-additions:3}
البروبان-2-أول؛ 2-ميثيل البروبان-2-أول؛ 1-ميثيل الهكسان الحلقي-1-أول (ماركوفنيكوف في
كل حالة).
\end{solution}

\begin{solution}{exo:b1:electrophilic-additions:4}
البروباين (\omterm{def:b1:electrophilic-additions:acetylide}{ألكاين طرفي}): \ce{CH3C#CH + NH2- -> CH3C#C- + NH3}. الإيثانول:
\ce{CH3CH2OH + NH2- -> CH3CH2O- + NH3}. أما البوت-2-اين والبروبين فلا هيدروجين
فيهما حمضيًا بما يكفي.
\end{solution}

\begin{solution}{exo:b1:electrophilic-additions:5}
يأخذ زوج $\pi$ الأيون \ce{H+} على كربون \ce{CH2}: الكاتيون الثالثي
\ce{(CH3)3C+}؛ ويرتبط به \omterm{def:b1:lewis-resonance-vsepr:lewis}{زوج حر} للماء: \ce{(CH3)3COH2+}؛ ويأخذ جزيء
ماء بروتونًا: \ce{(CH3)3COH}. وفي الحمض المركّز الساخن يكون الماء
شحيحًا ويفلت الألكين المتطاير: فيجري التفاعل العكسي، \omterm{def:b1:alcohol-activation:dehydration}{نزع الماء}.
\end{solution}

\begin{solution}{exo:b1:electrophilic-additions:6}
برتنة 2-ميثيل البروبين تعطي كاتيونًا ثالثيًا، وبرتنة البروبين كاتيونًا
ثانويًا. وللبرتنة المحددة للسرعة \omterm{def:b1:elementary-steps:energy-profile}{حالة انتقالية}
تشبه الكاتيون (هاموند): فيُبلغ الكاتيون الثالثي الأكثر ثباتًا
عبر حاجز أدنى.
\end{solution}

\begin{solution}{exo:b1:electrophilic-additions:7}
يجسر أيون البرومونيوم أحد وجهي الحلقة؛ ويفتحه البروميد من
الوجه الآخر: فتنتهي ذرتا البروم في وضع \textit{trans}. وثنائي البروميد \textit{trans}
\omterm{def:b1:stereochemistry-in-depth:stereocentre}{كيرالي} (لا مستوي مرآة فيه)، لكن البروميد يهاجم الكربونين
بالتساوي: فيتكوّن \omterm{def:b1:stereochemistry-in-depth:enantiomers}{المتماكبان المرآويان} بكميتين متساويتين، ويكون الناتج
راسيميًا، غير نشط ضوئيًا.
\end{solution}

\begin{solution}{exo:b1:electrophilic-additions:8}
برتنة C1 تعطي الكاتيون الثانوي على C2، بجوار C3 الثالثي
الحامل لهيدروجين؛ فينتقل ذلك الهيدروجين مع زوجه إلى C2 (انتقال
هيدريد 1,2)، تاركًا كاتيونًا ثالثيًا على C3، يأسره
\ce{Cl-}: 2-كلورو-2-ميثيل البوتان.
\end{solution}

\begin{solution}{exo:b1:electrophilic-additions:9}
البروباين + أميد الصوديوم: أيون البروبينيد؛ ومع 1-بروموبروبان (SN2):
\ce{CH3C#C-CH2CH2CH3}، أي هكس-2-اين. والبروبينيد مع البنزالدهيد، ثم حمض
مخفف: \ce{C6H5CH(OH)C#CCH3}، أي 1-فينيل بوت-2-اين-1-أول.
\end{solution}

\begin{solution}{exo:b1:electrophilic-additions:10}
البوت-2-ين-\cip{E}: برومونيوم على أحد الوجهين، وبروميد من الآخر على C2
أو على C3؛ والهجومان كلاهما يعطيان المركب نفسه، بالواصفين \cip{2R,3S}:
\omterm{def:b1:stereochemistry-in-depth:enantiomers}{مركب الميزو}، \omterm{def:b1:stereochemistry-in-depth:isomers}{متماكب فراغي} واحد. البوت-2-ين-\cip{Z}: الهجوم على C2 يعطي
\cip{2S,3S} (أو صورته المرآوية من الوجه الآخر)، والهجوم على C3
\cip{2R,3R}؛ بكميتين متساويتين: \omterm{def:b1:stereochemistry-in-depth:isomers}{متماكبان فراغيان}، \omterm{def:b1:stereochemistry-in-depth:enantiomers}{خليط راسيمي}.
\end{solution}

\begin{solution}{exo:b1:electrophilic-additions:11}
في أيون البرومونيوم غير المتناظر، يحمل الكربون الأكثر استبدالًا قسطًا أكبر
من الشحنة الموجبة (فرابطته C--Br أطول وأضعف، كأنه كاتيون يكاد يكون
ثالثيًا أو ثانويًا). فيهاجم الماء، الموجود بفائض كبير، ذلك الكربون
من الوجه المقابل للجسر: 1-بروموبروبان-2-أول.
\end{solution}

\begin{solution}{exo:b1:electrophilic-additions:12}
$D = 1$: ألكين. يعطي 2-ميثيل بوت-1-ين و2-ميثيل بوت-2-ين كلاهما
الكاتيون الثالثي، ومن ثم 2-برومو-2-ميثيل البوتان و2-ميثيل البوتان-2-أول.
ويختلف طيفا الرنين المغناطيسي النووي للبروتون فيهما: بروتونان فينيليان (\ce{=CH2}، أحاديتان) في
الأول، وبروتون واحد (\ce{=CH-}، رباعية) في الثاني.
\end{solution}

\begin{solution}{pb:b1:electrophilic-additions:1}
\textbf{1.} على كل كربون تتقدم \ce{CH3} على H: في \cip{Z} تقع مجموعتا الميثيل في
الجهة نفسها، وفي \cip{E} في جهتين متقابلتين.
\textbf{2.} \omterm{def:b1:stereochemistry-in-depth:isomers}{متماكبان فراغيان} ليس أحدهما صورة مرآوية للآخر: \omterm{def:b1:stereochemistry-in-depth:enantiomers}{متماكبان لامرآويان}.
\textbf{3.} الدوران حول C=C سيكسر الرابطة $\pi$، وهذا يكلّف طاقة أكبر
بكثير مما توفره الاصطدامات عند درجة حرارة الغرفة.
\textbf{4.} \cip{E}: لأن مجموعتي الميثيل فيه لا تتزاحمان.
\textbf{5.} \ce{C4H8 + Br2 -> C4H8Br2}.
\textbf{6.} مركزان فراغيان متكافئان: أربعة على الأكثر، وثلاثة في الواقع
(\cip{2R,3R} و\cip{2S,3S} والميزو \cip{2R,3S}).
\textbf{7.} يهاجم زوج $\pi$ إحدى ذرتي Br في \ce{Br2}، ويغادر زوج Br--Br
على هيئة \ce{Br-}، ويرتبط \omterm{def:b1:lewis-resonance-vsepr:lewis}{زوج حر} لذرة البروم المهاجَمة بالكربون
الثاني: أيون برومونيوم ثلاثي الحلقة.
\textbf{8.} في الأيون الجسري لكل ذرة ثمانية؛ أما \omterm{def:b1:electronic-effects:intermediates}{الكاتيون الكربوني} الثانوي
المفتوح ففيه كربون بستة إلكترونات.
\textbf{9.} يرتبط \omterm{def:b1:lewis-resonance-vsepr:lewis}{زوج حر} للأيون \ce{Br-} بأحد كربوني الحلقة؛ وتنكسر الرابطة C--Br
في الجسر، فيذهب زوجها إلى ذرة البروم الجسرية.
\textbf{10.} يشغل الجسر أحد الوجهين؛ وكما في SN2، يدخل \omterm{def:b1:electronic-effects:nucleophile}{المحب للنواة}
من الجهة المقابلة للرابطة التي تنكسر.
\textbf{11.} \omterm{def:b1:electrophilic-additions:halonium}{إضافة مضادة}.
\textbf{12.} كان الماء سينافس البروميد على أيون البرومونيوم ويعطي
بروموهيدرين.
\textbf{13.} \cip{2R,3S} (أو \cip{2S,3R}، وهو المركب نفسه).
\textbf{14.} المركب نفسه.
\textbf{15.} في \omterm{def:b1:stereochemistry-in-depth:isomers}{تشكّل} تكون فيه ذرتا البروم متطابقتين، يمر مستوي مرآة
بين C2 وC3؛ والواصفان متعاكسان.
\textbf{16.} \cip{2R,3R} و\cip{2S,3S}.
\textbf{17.} بكميتين متساويتين: \omterm{def:b1:stereochemistry-in-depth:enantiomers}{خليط راسيمي}.
\textbf{18.} نعم: الألكينان المتماكبان لامرآويًا يعطيان ناتجين مختلفين
(ميزو مقابل راسيمي).
\textbf{19.} لا دوران: يلغي \omterm{def:b1:stereochemistry-in-depth:enantiomers}{المتماكبان المرآويان} أحدهما الآخر.
\textbf{20.} $5.60/56.0 = \qty{0.100}{mol}$ لكل منهما.
\textbf{21.} $0.100 \times 0.85 \times 215.8 = \qty{18.3}{g}$ لكل منهما.
\textbf{22.} \textbf{$[\alpha] = 0^\circ$ (\omterm{def:b1:stereochemistry-in-depth:enantiomers}{مركب ميزو})، \qty{18.3}{g}}.
\end{solution}
